\subsection{Joint Distributions and Independence}
\begin{definition} \label{def:joint-distributions-and-marginals} \hlt{Joint and Marginal Distributions}
\begin{enumerate}[label=(\roman*)]
\setlength{\itemsep}{0pt}
\item The \hlt{joint mass function} of discrete $X,Y$ is $f_{X,Y}(x,y)\defeq\mathbb{P}(X=x,Y=y)$.
\item The \hlt{discrete marginals} are $f_X(x)=\sum\limits_y f_{X,Y}(x,y)$ and $f_Y(y)=\sum\limits_x f_{X,Y}(x,y)$.
\item The \hlt{continuous marginals} are $f_X(x)=\int_{-\infty}^{\infty}f_{X,Y}(x,y)\,dy$, and $f_Y(y)=\int_{-\infty}^{\infty}f_{X,Y}(x,y)\,dx$.
\item The \hlt{joint CDF} is $F_{X,Y}(x,y)\defeq\mathbb{P}(X\leq x,Y\leq y)$. $(X,Y)$ has joint density $f_{X,Y}$.
\end{enumerate}
\end{definition}

\begin{example} \label{ex:three-toss-joint-mass-table} \hlt{A Three-Toss Joint Mass Table}\\
Let $\Omega=\{H,T\}^3$ with $\mathbb{P}(\{\omega\})=2^{-3}$, set $X=\mathbf{1}_{\{\omega_1=H\}}$, $Y=\sum\limits_{i=1}^{3}\mathbf{1}_{\{\omega_i=H\}}$, and write $p_{xy}\defeq\mathbb{P}(X=x,Y=y)$.
\[
\begin{array}{c|cccc}
&Y=0&Y=1&Y=2&Y=3\\ \hline
X=0&\frac{1}{8}&\frac{2}{8}&\frac{1}{8}&0\\
X=1&0&\frac{1}{8}&\frac{2}{8}&\frac{1}{8}
\end{array}
\]
Thus $p_{0,2}=\frac{1}{8}$, $\sum\limits_y p_{0y}=\sum\limits_y p_{1y}=\frac{1}{2}$, and $\sum\limits_x p_{xj}=\binom{3}{j}2^{-3}$ for $j=0,\ldots,3$.
\end{example}

\begin{proposition} \label{prop:joint-cdf-rectangular-increments-and-density} \hlt{Rectangular Increments of a Joint CDF}\\
Let $F(x,y)\defeq\mathbb{P}(X\leq x,Y\leq y)$.
\begin{enumerate}[label=(\roman*)]
\setlength{\itemsep}{0pt}
\item If $x_1<x_2$ and $y_1<y_2$, then $\mathbb{P}(x_1<X\leq x_2,\ y_1<Y\leq y_2)=F(x_2,y_2)-F(x_1,y_2)-F(x_2,y_1)+F(x_1,y_1)$.
\item If $(X,Y)$ has density $f_{X,Y}$, $F\in C^2$, then $F(x,y)=\int_{-\infty}^{x}\int_{-\infty}^{y}f_{X,Y}(u,v)\,dv\,du$; $f_{X,Y}(x,y)=\partial_x\partial_yF(x,y)$.
\end{enumerate}
\begin{proof}
\begin{enumerate}[label=(\roman*)]
\setlength{\itemsep}{0pt}
\item Set $A_{ij}\defeq\{X\leq x_i,\ Y\leq y_j\}$ and $R\defeq\{x_1<X\leq x_2,\ y_1<Y\leq y_2\}$.\\
 Pointwise, $\mathbf{1}_R=\mathbf{1}_{A_{22}}-\mathbf{1}_{A_{12}}-\mathbf{1}_{A_{21}}+\mathbf{1}_{A_{11}}$; applying $\mathbb E$ gives the identity.
\item The density assumption gives $F(x,y)=\mathbb P(X\leq x,Y\leq y)=\int_{-\infty}^{x}\int_{-\infty}^{y}f_{X,Y}(u,v)\,dv\,du$.\\
Since $F\in C^2$, the two-variable FTC gives $\partial_x\partial_yF=f_{X,Y}$.
\end{enumerate}
\hlt{Structural Commentary}: Joint CDFs encode southwest rectangles; finite rectangles are four-corner increments.\\
\hlt{Key Insight}: Mixed derivatives are infinitesimal rectangular increments.
\end{proof}
\end{proposition}

\begin{figure}[H]
  \centering
  % Source integration: schematic redraw of the rectangular increment formula for
% a bivariate CDF from the attached joint-distribution chapter screenshots.
\begin{tikzpicture}[
  x=1.05cm,
  y=0.85cm,
  >=Stealth,
  line cap=round,
  line join=round,
  every node/.style={font=\scriptsize},
  axis/.style={->, gray!70, line width=0.45pt},
  tick/.style={gray!70, line width=0.4pt},
  guide/.style={gray!40, dashed, line width=0.4pt},
  event/.style={fill=blue!14, draw=blue!70!black, line width=0.85pt},
  corner/.style={circle, fill=blue!70!black, inner sep=0pt, minimum size=2.4pt},
  axis label/.style={fill=white, inner sep=1pt, text=black!78},
  event label/.style={font=\small, text=blue!70!black}
]
  \coordinate (lower left) at (1.25,1.05);
  \coordinate (lower right) at (3.25,1.05);
  \coordinate (upper left) at (1.25,2.65);
  \coordinate (upper right) at (3.25,2.65);

  \draw[axis] (-0.35,0) -- (4.15,0) node[black!78,right] {$x$};
  \draw[axis] (0,-0.3) -- (0,3.35) node[black!78,above] {$y$};

  \draw[guide] (1.25,0) -- (lower left);
  \draw[guide] (3.25,0) -- (lower right);
  \draw[guide] (0,1.05) -- (lower left);
  \draw[guide] (0,2.65) -- (upper left);

  \path[event] (lower left) rectangle (upper right);
  \foreach \p in {lower left,lower right,upper left,upper right}
    \node[corner] at (\p) {};

  \draw[tick] (1.25,0.025) -- (1.25,-0.025) node[axis label, below=2pt] {$x_1$};
  \draw[tick] (3.25,0.025) -- (3.25,-0.025) node[axis label, below=2pt] {$x_2$};
  \draw[tick] (0.025,1.05) -- (-0.025,1.05) node[axis label, left=2pt] {$y_1$};
  \draw[tick] (0.025,2.65) -- (-0.025,2.65) node[axis label, left=2pt] {$y_2$};

  \node[event label] at (2.25,1.85) {$R$};
\end{tikzpicture}

  \caption{The rectangular probability is a four-term increment of the joint CDF.}
  \label{fig:joint-cdf-rectangular-increment}
\end{figure}

\begin{example} \label{ex:bivariate-density-nonrectangular-region} \hlt{A Non-Rectangular Bivariate Density}\\
Let $S=\{(x,y):x^2\leq y\leq1\}$ and $f_{X,Y}(x,y)=cx^2y\mathbf{1}_S(x,y)$.\\
Since $S=\{(x,y):-1\leq x\leq1,\ x^2\leq y\leq1\}$, $1=c\int_{-1}^{1}\int_{x^2}^{1}x^2y\,dy\,dx=\frac{4c}{21}$, so $c=\frac{21}{4}$.\\
For $E=\{X\geq Y\}$, $S\cap E=\{(x,y):0\leq x\leq1,\ x^2\leq y\leq x\}$, hence $\mathbb{P}(E)=(\frac{21}{4})\int_0^1\int_{x^2}^{x}x^2y\,dy\,dx=\frac{3}{20}$.\\
\hlt{Key Insight}: In bivariate integration, pick one variable first and let the other range over the corresponding slice.
\end{example}

\begin{figure}[H]
  \centering
  % Source integration: simplified redraw of the non-rectangular bivariate-density
% support and the event region from the attached random-variable chapter screenshots.
\begin{tikzpicture}[
  x=3.0cm,
  y=3.0cm,
  >=Stealth,
  line cap=round,
  line join=round,
  every node/.style={font=\scriptsize},
  axis/.style={->, gray!70, line width=0.45pt},
  tick/.style={gray!70, line width=0.4pt},
  support/.style={fill=gray!12, draw=none},
  event/.style={fill=blue!22, draw=none},
  boundary/.style={black!72, line width=0.75pt},
  event boundary/.style={blue!70!black, line width=0.9pt}
]
  \path[support, domain=-1:1, variable=\x]
    plot ({\x},{\x*\x}) -- (1,1) -- (-1,1) -- cycle;
  \path[event] (0,0) -- plot[domain=0:1, variable=\x] ({\x},{\x*\x}) -- cycle;

  \draw[axis] (-1.18,0) -- (1.22,0) node[black!78,right] {$x$};
  \draw[axis] (0,-0.08) -- (0,1.18) node[black!78,above] {$y$};

  \draw[boundary, domain=-1.04:1.04, samples=120, variable=\x] plot ({\x},{\x*\x});
  \draw[boundary] (-1,1) -- (1,1);
  \draw[event boundary] (0,0) -- (1,1);

  \foreach \x/\lab in {-1/-1,0/0,1/1}
    \draw[tick] (\x,0.018) -- (\x,-0.018) node[black!78,below] {$\lab$};
  \node[black!78,anchor=south east] at (-0.05,1.02) {$1$};
\end{tikzpicture}

  \caption{\textcolor{gray!70!black}{Gray}: support $x^2\leq y\leq1$; \textcolor{blue!70!black}{blue}: event $X\geq Y$ ($x^2\leq y\leq x$).}
  \label{fig:parabolic-bivariate-region}
\end{figure}

\begin{definition} \label{def:independent-random-variables-and-iid-samples} \hlt{Independent Random Variables and IID Samples}
\begin{enumerate}[label=(\roman*)]
\setlength{\itemsep}{0pt}
\item Random variables $X_1,\ldots,X_n$ are \hlt{independent} if $\mathbb{P}(X_1\in B_1,\ldots,X_n\in B_n)=\prod\limits_{i=1}^{n}\mathbb{P}(X_i\in B_i)$ $\forall B_i\in\mathcal{B}(\R)$.
\item They are \hlt{identically distributed} if $\mathcal{L}(X_1)=\cdots=\mathcal{L}(X_n)$.
\item They are \hlt{iid} if they are independent and identically distributed.
\item The \hlt{random vector} $X=(X_1,\ldots,X_n)$ is a \hlt{random sample} when $X_1,\ldots,X_n$ are iid from a common law.
\end{enumerate}
\end{definition}

\begin{proposition} \label{prop:factorization-tests-for-random-variable-independence} \hlt{Factorization Tests for Independence}
\begin{enumerate}[label=(\roman*)]
\setlength{\itemsep}{0pt}
\item \hlt{Discrete factorization}: discrete $X,Y$ are independent iff $f_{X,Y}(x,y)=f_X(x)f_Y(y)$ for all $x,y$.
\item \hlt{Density factorization}: if $(X,Y)$ has joint density $f_{X,Y}$ and marginals $f_X,f_Y$,\\
then $X,Y$ are independent iff $f_{X,Y}(x,y)=f_X(x)f_Y(y)$ a.e.
\item \hlt{Rectangular support test}: if $f_{X,Y}(x,y)=g(x)h(y)\mathbf{1}_{A\times B}(x,y)$ for measurable $A,B$,\\
then $X,Y$ are independent after normalization.
\end{enumerate}
\begin{proof}
\begin{enumerate}[label=(\roman*)]
\setlength{\itemsep}{0pt}
\item[$(i)$] ($\Rightarrow$) If $X\perp\!\!\!\perp Y$, then $f_{X,Y}(x,y)=\mathbb{P}(X=x,Y=y)=\mathbb{P}(X=x)\mathbb{P}(Y=y)=f_X(x)f_Y(y)$.
\item[$(i)$] ($\Leftarrow$) If $f_{X,Y}=f_Xf_Y$, then for $A,B$, $\mathbb{P}(X\in A,Y\in B)=\sum\limits_{x\in A}\sum\limits_{y\in B}f_X(x)f_Y(y)=\mathbb{P}(X\in A)\mathbb{P}(Y\in B)$.
\item[$(ii)$] ($\Rightarrow$) If $X\perp\!\!\!\perp Y$, the product law has density $f_Xf_Y$, so uniqueness of densities gives $f_{X,Y}=f_Xf_Y$ a.e.
\item[$(ii)$] ($\Leftarrow$) If $f_{X,Y}=f_Xf_Y$ a.e., then for Borel $A,B$,\\
$\mathbb{P}(X\in A,Y\in B)=\int_A\int_Bf_X(x)f_Y(y)\,dy\,dx=\mathbb{P}(X\in A)\mathbb{P}(Y\in B)$.
\item[$(iii)$] Let $G=\int_Ag$ and $H=\int_Bh$. Normalization gives $GH=1$, while $f_X(x)=g(x)H\mathbf{1}_A(x)$ and\\
$f_Y(y)=h(y)G\mathbf{1}_B(y)$, hence $f_X(x)f_Y(y)=f_{X,Y}(x,y)$ a.e.; apply (ii).
\end{enumerate}
\hlt{Structural Commentary}: Independence is product structure of the joint law.\\
\hlt{Key Insight}: Independence means the joint law factors.
\end{proof}
\end{proposition}
